%---------------------------------------------------------------------------------------------------
% Einstellungen
% (gelten nur in Zusammenarbeit mit pdflatex)
%---------------------------------------------------------------------------------------------------
\documentclass[
  pagesize, % flexible Auswahl des Papierformats
  a4paper, % DIN A4
  oneside, % einseitiger Druck
  BCOR5mm, % Bindungskorrektur
  headsepline, % Strich unter der Kopfzeile
  11pt, % 11pt Schriftgröße
  halfparskip, % Europäischer Satz: Abstand zwischen Absätzen
  abstracton, % Spezielle Formatierung, die erlaubt, dass die Zusammenfassung vor dem Inhaltsverzeichnis steht
  %draft, % Es handelt sich um eine Vorabversion
  final, % Es handelt sich um die endgültige Version
  liststotoc, % Tabellen- und Abbildungsverzeichnis im Inhaltsverzeichnis
  idxtotoc, % Index im Inhaltsverzeichnis
  bibtotoc,                                            % Literaturverzeichnis im Inhaltsverzeichnis
]{scrartcl}                                            % KOMA-Scriptklasse Report

%---------------------------------------------------------------------------------------------------
\usepackage[english,ngerman]{babel}                    % deutsche Trennmuster
\usepackage[T1]{fontenc}                               % EC-Schriften, Trennstellen nach Umlauten
\usepackage[utf8x]{inputenc}                          % direkte Umlauteingabe (ä statt "a)
                                                       % latin1/latin9 für unixoide Systeme
                                                       % (latin1 ist auch unter Win verwendbar)
                                                       % ansinew für Windows
                                                       % applemac Macs
                                                       % cp850 OS/2
\usepackage{times}                                     % Schriften Paket
\usepackage{array,ragged2e}                            % Wichtig für Abstandsformatierung

%---------------------------------------------------------------------------------------------------
\usepackage{cmbright}                                  % serifenlose Schrift als Standard
                                                       % + alle für TeX benötigten mathematischen
                                                       %   Schriften einschließlich der AMS-Symbole
\usepackage[scaled=.90]{helvet}                        % skalierte Helvetica als \sfdefault
\usepackage{courier}                                   % Courier als \ttdefault

%---------------------------------------------------------------------------------------------------
\usepackage[automark]{scrpage2}                        % Anpassung der Kopf- und Fußzeilen
\usepackage{xspace}                                    % Korrekter Leerraum nach Befehlsdefinitionen
\usepackage{setspace}                                  % Dieses Package brauchen wir für den anderthalbzeiligen Abstand.
\usepackage[numbers]{natbib}                           % Neuimplementierung des \cite-Kommandos
%\usepackage{bibgerm}                                  % Deutsche Bezeichnungen (deprecated)
\usepackage[german]{babelbib}                          % Deutsche Bezeichnungen
\usepackage[absolute, overlay]{textpos}                % placing boxes at absolute positions
\usepackage[final]{pdfpages}                           % include pages of external PDF documents
\usepackage{tabularx}                                  % Spaltenbreite bis zur Seitenbreite dehnen
\usepackage{makeidx}                                   % Paket zur Erstellung eines Stichwortverzeichnisses
\makeindex                                             % Automatische Erstellung des Stichwortverzeichnis
\usepackage[intoc,german,prefix]{nomencl}

% \usepackage{footnote} % Ermöglicht Fußnoten in gleitenden Umgebungen
% \usepackage{capt-of} % Ermöglicht Beschriftungen beliebiger Objekte

\makenomenclature

%---------------------------------------------------------------------------------------------------
 \usepackage{graphicx}                                 % Zur Einbindung von PDF-Bildern
 \usepackage[colorlinks,                               % Einstellen und Laden des Hyperref-Pakets
  pdftex,
  bookmarks,
  bookmarksopen=false,
  bookmarksnumbered,
  citecolor=blue,
  linkcolor=blue,
  urlcolor=blue,
  filecolor=blue,
  linktocpage,
  pdfstartview=Fit,                                  % startet mit Ganzseitenanzeige
  pdfsubject={Hausarbeit im Rahmen der Vorlesung Modellierung technischer Systeme},
  pdftitle={Hausarbeit im Rahmen der Vorlesung Modellierung technischer Systeme},
  pdfauthor={Andreas Krohn}]{hyperref}
% \pdfcompresslevel=9
%\usepackage{wrapfig}

%---------------------------------------------------------------------------------------------------
% Inhaltsverzeichnis und Abschnittnummerierung
%---------------------------------------------------------------------------------------------------
\setcounter{secnumdepth}{2}   % Ich habe recht kurze Kapitel. Die sollen nicht durchnummeriert sein.
\setcounter{tocdepth}{2}

%---------------------------------------------------------------------------------------------------
% Abbildungsverzeichnis
%---------------------------------------------------------------------------------------------------
\graphicspath{{graphics/}}

%---------------------------------------------------------------------------------------------------
% Kopf- und Fußzeilen
%---------------------------------------------------------------------------------------------------
\pagestyle{scrheadings}
\clearscrheadings
\clearscrplain
\clearscrheadfoot
\ohead{\pagemark}
\ihead{\headmark}

%---------------------------------------------------------------------------------------------------
% Neue Befehle
%---------------------------------------------------------------------------------------------------
%---------------------------------------------------------------------------------------------------
% Neue Befehle
%---------------------------------------------------------------------------------------------------

%---------------------------------------------------------------------------------------------------
% Umbenennen des Symbolverzeichnisses
%---------------------------------------------------------------------------------------------------
\renewcommand{\nomname}{Glossar}				% Das Symbolverzeichnis heisst nun "Glossar"
\renewcommand{\nomlabel}[1]{						% Die zu erklärenden Begriffe sind nun fett hervorgehoben
	\hfil \textbf{#1} \hfil
}

%---------------------------------------------------------------------------------------------------
% Ein paar ganz nützliche Befehle von Lars Mählmann
%---------------------------------------------------------------------------------------------------
%für Kommentare
\newcommand{\colb}{\color{green}}
\newcommand{\colbl}{\color{black}}

%---------------------------------------------------------------------------------------------------
% Befehle zum Erstellen des Index
% \addIndexEntry{Eintrag in den Index}
% \addSubIndexEntry{Eintrag in den Index}{Eintrag des übergeordneten Eintrags}
%---------------------------------------------------------------------------------------------------
\newcommand{\addIndexEntry}[1]{#1\index{#1}}
\newcommand{\addSubIndexEntry}[2]{#1\index{#2!#1}}

%---------------------------------------------------------------------------------------------------
% LaTeX in eigenem Font
%---------------------------------------------------------------------------------------------------
\newcommand{\myLatex}{
	{\rmfamily\LaTeX\xspace}
}

%---------------------------------------------------------------------------------------------------
% Befehl zum Erstellen und Hervorheben eines Zitats
% Parameter:
% 1. Zitat
% 2. Author
% 3. Quelle
%---------------------------------------------------------------------------------------------------
\newcommand{\myCitation}[3]{
	\begin{flushright}
	\begin{minipage}{.4\linewidth}
		\footnotesize\rmfamily\itshape
		#1 \\
		\RaggedLeft #2 \\
		#3
	\end{minipage}
	\end{flushright}
	\nobreakspace
}

%---------------------------------------------------------------------------------------------------
% Erstellung von Deckblatt (Seite 1) und Titelblatt (Seite 2)
%---------------------------------------------------------------------------------------------------
\newcommand{\createCoverAndTitlePage}[8]{
	\createCover{#1}{#2}{#3}
	\createTitlePage{#1}{#2}{#3}{#4}{#5}{#6}{#7}{#8}
}

%---------------------------------------------------------------------------------------------------
% Erstellung von Deckblatt (Seite 1)
% Anwendung:
% \createCover{Art der Arbeit}{Autor}{Titel}
%---------------------------------------------------------------------------------------------------
\newcommand{\createCover}[3]{
	\thispagestyle{empty}
	\begin{titlepage}

	\setlength{\TPHorizModule}{1mm}
	\setlength{\TPVertModule}{1mm}
	\textblockorigin{0mm}{0mm} % start everything near the top-left corner

	% Art der Arbeit
	\begin{textblock}{111}(83,115)
		\begin{minipage}[c][1,78cm][c]{11,09cm}
  		\fontsize{22pt}{20pt}
  		\selectfont
  		\begin{center}
  		#1
  		\end{center}
		\end{minipage}
	\end{textblock}

	% Name & Titel
	\begin{textblock}{111}(83,131)
		\begin{minipage}[c][4,81cm][t]{11,09cm}
		\linespread{1.2}
    \fontsize{16pt}{14pt}
    \selectfont
    \begin{center}
    #2 \\ \medskip
    #3
    \end{center}
    \end{minipage}
	\end{textblock}
	% Infos zur Arbeit und zum Fachbereich
	\begin{textblock}{186}(22,264)
  	\begin{minipage}[t][5,72cm][l]{17,57cm}
    	\fontsize{12pt}{12pt}
    	\selectfont
			{\em Fakultät Technik und Informatik \hfill Faculty of Engineering and Computer Science}\\
			{\em Department Informatik \hfill Department of Computer Science}
  	\end{minipage}
	\end{textblock}
	\newpage
	\end{titlepage}
%---------------------------------------------------------------------------------------------------
% Wichtig! Entsprechendes Auskommentieren!
%---------------------------------------------------------------------------------------------------
	\includepdf{pdf/titel}           				% zum Ausdruck auf blanko Papier
  %\includepdf{pdf/titel_leer}           	% zum Ausdruck auf die Pappe
}

%---------------------------------------------------------------------------------------------------
% Erstellung von Titelblatt (Seite 2)
% Anwendung:
% \createTitlePage{Art der Arbeit}{Author}{Titel}{Studiengang}{Erstprüfer}{Zweitprüfer}
%---------------------------------------------------------------------------------------------------
\newcommand{\createTitlePage}[7]{
	\thispagestyle{empty}

	\setlength{\TPHorizModule}{1mm}
	\setlength{\TPVertModule}{\TPHorizModule}
	\textblockorigin{0mm}{0mm} % start everything near the top-left corner

	% Name & Titel
	\begin{textblock}{130}(40,63)
		\begin{minipage}[c][5,9cm][t]{13cm}
			\begin{center}
			\linespread{1.2}
			\fontsize{18pt}{18pt}
  		\selectfont
  		#2 \\ \medskip
  		\fontsize{16pt}{16pt}
  		#3
  		\end{center}
		\end{minipage}
	\end{textblock}

	% Infos zur Arbeit und zum Fachbereich
	\begin{textblock}{126}(32,214)
  	\begin{minipage}[t][5,72cm][l]{12,57cm}
    	\fontsize{12pt}{12pt}
    	\selectfont
    	#1 eingereicht im Rahmen der #1prüfung\\
    	im Studiengang #4\\
			am Department Informatik\\
			der Fakultät Technik und Informatik\\
			der Hochschule für Angewandte Wissenschaften Hamburg\\
			\\\
			\\\
			Betreuender Prüfer : #5\\
			Zweitgutachter : #6\\
			\\\
			Abgegeben am \today
  	\end{minipage}
	\end{textblock}
	\	% WICHTIG! Damit wird nach dem Titelblatt eine neue Seite angefangen! Sonst werden Titelblatt &
  	% Danksagung auf eine Seite gedruckt!
}

%---------------------------------------------------------------------------------------------------
% Versicherung über Selbstständigkeit
%---------------------------------------------------------------------------------------------------
\newcommand{\asurency}{
	\chapter*{Versicherung über Selbstständigkeit}
	\vfill
	Hiermit versichere ich, dass ich die vorliegende Arbeit im Sinne der Prü\-fungs\-ord\-nung nach \S 24(5) 		ohne fremde Hilfe selbstständig verfasst und nur die angegebenen Hilfsmittel benutzt habe.
	\vfill
	\begin{tabularx}{\linewidth}{X l X}
	Hamburg, \today	& \qquad \qquad \qquad	& \\
	\cline{1-1}
	\cline{3-3}
	Ort, Datum	& \qquad \qquad \qquad	& Unterschrift \\
	\end{tabularx}
	\vfill
	\vfill
	\vfill
}

%---------------------------------------------------------------------------------------------------
% Fügt ein Wort dem Index zu
%---------------------------------------------------------------------------------------------------
\newcommand{\toIndex}[1]{#1\index{#1}}

%---------------------------------------------------------------------------------------------------
% Dient zum Eintragen folgender Dinge in die Zusammenfassung (Abstract):
%	- Thema
% - Stichworte
% - Kurzfassung
% Benutzung wie folgt:
% \abstractentry{Titel}{Text}
%---------------------------------------------------------------------------------------------------
\newcommand{\abstractentry}[2]{
	\textbf{\large#1}\\
	\nobreakspace
	\begin{tabular}{lp{142mm}}
		\hspace*{7mm} & #2 \\
	\end{tabular}
	\vfill
}

%---------------------------------------------------------------------------------------------------
% Erstellt eine Defintion
% Anwendung: \definition{Die Definition}
%---------------------------------------------------------------------------------------------------
\newcommand{\definition}[1]{
\begin{tabular}[ht]{lp{135mm}}
	\textbf{Def.:} & #1 \\
\end{tabular}
}

%---------------------------------------------------------------------------------------------------
% Erstellt eine Widmung
% Anwendung: \dedication{Wem ist das Schriftstück gewidmet}
%---------------------------------------------------------------------------------------------------
\newcommand{\createDedication}[1]{
	\newpage
	\thispagestyle{empty}
	\begin{tabular}{lp{60mm}}
		\hspace*{100mm} & \itshape\rmfamily#1 \\
	\end{tabular}
	\vfill
}

%---------------------------------------------------------------------------------------------------
% Häufig verwendete Namen mit Literaturverweis und Indexeintrag
%--------------------------------------------------------------------------------------------------
\newcommand{\butrynowski}{Christian Butrynowski\index{Butrynowski, Christian} \citep{Butrynowski:2005}\xspace}
\newcommand{\luepke}{André Lüpke\index{Lüpke, André}\citep{Luepke:2004}\xspace}
\newcommand{\bresch}{Marco Bresch\index{Bresch, Marco} \citep{Bresch:2004}\xspace}


%---------------------------------------------------------------------------------------------------
% Die folgenden Befehle wurden aus der Vorlage von Michael Knop übernommen
%--------------------------------------------------------------------------------------------------
%---------------------------------------------------------------------------------------------------
% Ident
%---------------------------------------------------------------------------------------------------
\newcommand{\ident}[1]{                             % ein Parameter
	\small\ttfamily#1\sffamily\normalsize
}

%---------------------------------------------------------------------------------------------------
% Kürzel
%---------------------------------------------------------------------------------------------------
% Hier sind Makros definiert, die die Eingabe erleichtern sollen. Für korrekte Abstände zwischen
% "z.B." sorgt also ein "z.\,B." (LaTeX-Befehl für kleineren Abstand)
% Schneller schreibt sich das durch das Makro "\zB":

% \newcommand{\zB}{z.\,B.\ }

% Hier ist der Rest aber mit dem Paket xspace verwirklicht. Damit kann
% man bei Bedarf den Abstand mit "\hspace" exakt eingeben. Dann zeigt
% LaTeX keine Toleranz bei den Abkürzungen und macht eben exakt das
% untenstehende.

%\renewcommand{\entryname}{K\"urzel}
%\renewcommand{\descriptionname}{Beschreibung}

\newcommand{\vgl}{vgl.\@\xspace}
\newcommand{\abb}{Abb.\@\xspace}
\newcommand{\zB}{z.\nolinebreak[4]\hspace{0.125em}\nolinebreak[4]B.\@\xspace}
\newcommand{\bzw}{bzw.\@\xspace}
\newcommand{\dahe}{d.\nolinebreak[4]\hspace{0.125em}h.\nolinebreak[4]\@\xspace}
\newcommand{\etc}{etc.\@\xspace}
\newcommand{\bzgl}{bzgl.\@\xspace}
\newcommand{\so}{s.\nolinebreak[4]\hspace{0.125em}\nolinebreak[4]o.\@\xspace}
\newcommand{\iA}{i.\nolinebreak[4]\hspace{0.125em}\nolinebreak[4]A.\@\xspace}
\newcommand{\sa}{s.\nolinebreak[4]\hspace{0.125em}\nolinebreak[4]a.\@\xspace}
\newcommand{\su}{s.\nolinebreak[4]\hspace{0.125em}\nolinebreak[4]u.\@\xspace}
\newcommand{\ua}{u.\nolinebreak[4]\hspace{0.125em}\nolinebreak[4]a.\@\xspace}
\newcommand{\og}{o.\nolinebreak[4]\hspace{0.125em}\nolinebreak[4]g.\@\xspace}

\newcommand{\HAW}{Hochschule für Angewandte Wissenschaften Hamburg\xspace}
\newcommand{\GNU}{GNU\xspace}
\newcommand{\GPL}{\GNU Public License\xspace}

\newcommand{\ACM}{ACM\xspace}
\newcommand{\PDA}{PDA\xspace}


%---------------------------------------------------------------------------------------------------
% Trennung
%---------------------------------------------------------------------------------------------------
%---------------------------------------------------------------------------------------------------
% Trennung
% Hier können alle Wörtertrennungen definiert werden. Die nachfolgenden dienen als Beispiel
% und wurden aus der Vorlage von Michael Knop übernommen.
%---------------------------------------------------------------------------------------------------
\hyphenation{Web-ap-pli-ka-tion Web-ap-pli-ka-tio-nen Web-an-wen-dung Web-an-wen-dung-en My-SQL Kon-text-in-for-ma-ti-onen}

%---------------------------------------------------------------------------------------------------
% Anpassung der Parameter, die TeX bei der Berechnung der Zeilenumbrüche verwendet:
%---------------------------------------------------------------------------------------------------
\tolerance 1414
\hbadness 1414
\emergencystretch 1.5em
\hfuzz 0.3pt
\widowpenalty=10000
\vfuzz \hfuzz
\raggedbottom
